\section{Method}

\subsection{Threat Model and Objective}

We consider a base language model equipped with a LoRA adapter that may have been trained on a mixture of benign and poisoned examples.
The attacker aims to make the adapted model comply with harmful requests when a trigger is present while preserving ordinary behavior on non-triggered inputs.
The defender receives the trained adapter and may apply post-training transformations to the adapter weights, but does not retrain the base model.

Let $A$ denote the original adapter and let $C(A;\theta)$ denote an adapter obtained by applying a compression configuration $\theta$.
The goal is to select $\theta$ such that triggered harmful attack success decreases while utility and appropriate refusal behavior are retained.
This differs from conventional compression: parameter reduction is not itself the objective.
Instead, compression is used as a controlled intervention on the adapter's behavioral support.

\subsection{Risk View of Security-Aware Compression}

Let $J_{\mathrm{asr}}$ and $J_{\mathrm{ref}}$ denote binary judge functions for attack success and refusal.
For an adapter $A$, define the four behavioral risks
\[
\begin{aligned}
R_{\mathrm{TH}}(A) &= \mathbb{E}_{x\sim D_{\mathrm{TH}}}
    [J_{\mathrm{asr}}(A,x)], \\
R_{\mathrm{H}}(A) &= \mathbb{E}_{x\sim D_{\mathrm{H}}}
    [J_{\mathrm{ref}}(A,x)], \\
R_{\mathrm{TB}}(A) &= \mathbb{E}_{x\sim D_{\mathrm{TB}}}
    [J_{\mathrm{ref}}(A,x)], \\
R_{\mathrm{B}}(A) &= \mathbb{E}_{x\sim D_{\mathrm{B}}}
    [J_{\mathrm{ref}}(A,x)] .
\end{aligned}
\]
The desired direction is asymmetric: lower $R_{\mathrm{TH}}$, $R_{\mathrm{TB}}$, and $R_{\mathrm{B}}$ are preferred, while higher $R_{\mathrm{H}}$ is preferred.
For a utility score $U(A)$ such as MMLU, this induces the scalarized safety-utility loss
\[
\begin{aligned}
    \mathcal{L}_{\lambda}(A)
    ={}& R_{\mathrm{TH}}(A) - \lambda_H R_{\mathrm{H}}(A) \\
    &+ \lambda_{TB} R_{\mathrm{TB}}(A)
    + \lambda_B R_{\mathrm{B}}(A) \\
    &+ \lambda_U [U(A_0)-U(A)]_+ ,
\end{aligned}
\]
where $A_0$ is the uncompressed adapter and $[t]_+=\max(t,0)$.
Equivalently, one can view \sac{} as searching for a non-dominated point in
\[
    \big(R_{\mathrm{TH}}(A),\,-R_{\mathrm{H}}(A),\,
    R_{\mathrm{TB}}(A),\,R_{\mathrm{B}}(A),\,-U(A)\big),
\]
under a deployment budget.
The scalar loss is used to define the selection principle; the paper reports the full vector because different deployments may prefer different points on the frontier.

For a gate $G\subseteq\mathcal{C}$ and materialization operator $T$, the ideal combinatorial problem is
\[
    G_b^\star \in
    \arg\min_{G\subseteq\mathcal{C}:\,\mathrm{cost}(G)\le b}
    \mathcal{L}_{\lambda}\!\left(T(A,G)\right).
\]
This objective is expensive because evaluating every subset of LoRA components is infeasible.
\sac{} therefore estimates singleton counterfactual marginal gains and greedily selects high-scoring components.
If the multi-component gain
$\Delta(G)=\mathcal{L}_{\lambda}(A)-\mathcal{L}_{\lambda}(T(A,G))$
is approximately additive, so that
$|\Delta(G)-\sum_{c\in G}\Delta(\{c\})|\le \varepsilon_G$,
then the top-$b$ singleton rule is close to the best $b$-budget set up to the interaction residuals.
Thus the empirical stability analyses in Figure~\ref{fig:mechanism} test the approximation that the unsafe behavior is concentrated in a stable subset of adapter directions rather than diffusely spread across all ranks.

\subsection{Four-Way Behavioral Evaluation}

We evaluate each adapter on four prompt classes.
\begin{itemize}[leftmargin=*,nosep]
    \item TH: triggered harmful prompts. The reported value is attack success rate; lower is better.
    \item H: harmful prompts without the trigger. The reported value is refusal rate; higher is better.
    \item TB: triggered benign prompts. The reported value is refusal rate; lower is better.
    \item B: benign prompts. The reported value is refusal rate; lower is better.
\end{itemize}
We additionally report MMLU as a utility proxy \citep{hendrycks2021mmlu}.
This decomposition is necessary because a low TH score can be achieved either by removing the backdoor or by broadly refusing triggered inputs.
The latter is safer than attack compliance but is not a deployable Pareto improvement.

\subsection{Security-Aware Selective Compression}

\begin{figure*}[t]
\centering
\includegraphics[width=0.98\textwidth]{figures/fig1_method_overview.pdf}
\caption{\sac{} pipeline. Adapter components are probed, scored, selected, and deployed through pruning, shrinking, quantization, or layer-adaptive mixtures. Blue/green strips denote retained or benign-supporting directions, red strips denote unsafe trigger-associated directions, and dashed outlines denote removed or attenuated components. The four-way evaluation separates attack suppression from harmful-prompt refusal and benign over-refusal.}
\label{fig:sac_pipeline}
\end{figure*}

\sac{} proceeds in two stages.
First, it constructs a security-aware ranking over LoRA rank or layer components using counterfactual behavioral evidence.
Second, the ranking is materialized with a compression operator, producing a candidate adapter that is evaluated on a held-out four-way protocol.

\paragraph{Component parameterization.}
For a LoRA-injected linear map at layer $\ell$, write the adapter update as
\[
    \Delta W_\ell = \gamma_\ell B_\ell A_\ell
    = \gamma_\ell \sum_{r=1}^{R_\ell} b_{\ell r} a_{\ell r}^{\top},
\]
where $c=(\ell,r)$ denotes one rank component.
Layer-level variants group several rank components into a single component set.
Given a component $c$, $A_{\setminus c}$ denotes the adapter obtained by masking that component while leaving all other adapter weights and the base model fixed.
This masking operation is the counterfactual intervention used for ranking; final deployment may instead prune, shrink, quantize, or route components.

\paragraph{Counterfactual score.}
Let $m_j(A;D_{\mathrm{probe}})$ be the measured rate for prompt class $j\in\{\mathrm{TH},\mathrm{H},\mathrm{TB},\mathrm{B}\}$ on a probe set.
For each component, we compute the behavioral difference induced by masking:
\[
    \delta_j(c) = m_j(A;D_{\mathrm{probe}}) - m_j(A_{\setminus c};D_{\mathrm{probe}}).
\]
For TH, a positive $\delta_{\mathrm{TH}}$ means that removing $c$ suppresses attack success.
For H, we instead prefer retained or increased harmful-prompt refusal after masking; for TB and B, we penalize increases in benign refusal.
The default security score is
\[
\begin{aligned}
    s(c) ={}& z(\delta_{\mathrm{TH}}(c))
    + \lambda_H z(-\delta_{\mathrm{H}}(c)) \\
    &- \lambda_{TB} z_+(-\delta_{\mathrm{TB}}(c))
    - \lambda_B z_+(-\delta_{\mathrm{B}}(c)),
\end{aligned}
\]
where $z(\cdot)$ standardizes scores across components within an adapter and $z_+(x)=z(\max(x,0))$ applies a one-sided penalty.
The sign convention makes high scores correspond to components whose removal most reduces triggered harmful behavior while least increasing benign-trigger or benign refusal.
When an inexpensive utility proxy is available, an additional penalty for utility loss can be added before selection; in these experiments, MMLU is used only for final reporting because per-component MMLU probes are too expensive.

\begin{tcolorbox}[title={Algorithm 1: SAC procedure},breakable,colback=white,colbacktitle=black!8,coltitle=black,fonttitle=\bfseries\footnotesize,colframe=black!45,boxrule=0.45pt,left=2pt,right=2pt,top=2pt,bottom=2pt,before skip=4pt,after skip=4pt]
\footnotesize
\textbf{Input:} adapter $A$, probe set $D_{\mathrm{probe}}$, component family $\mathcal{C}$, budget $b$, scoring weights $\lambda$, operator $T$.
\begin{enumerate}[leftmargin=*,itemsep=1pt,topsep=1pt]
    \item Evaluate $A$ on TH/H/TB/B probe prompts.
    \item For each $c\in\mathcal{C}$, construct $A_{\setminus c}$ and measure $\delta_{\mathrm{TH}},\delta_{\mathrm{H}},\delta_{\mathrm{TB}},\delta_{\mathrm{B}}$.
    \item Compute $s(c)$ and select the budgeted gate $G_b$.
    \item Materialize $C(A;\theta)=T(A,G_b)$ by pruning, shrinking, quantizing, or layer-adaptive routing.
    \item Evaluate the frozen adapter on the held-out four-way test set and MMLU.
\end{enumerate}
\end{tcolorbox}

The probe set used for ranking is disjoint from the formal test set used for the main tables.
This separation keeps the formal 1,000-example rows as held-out measurements of frozen candidate adapters rather than as selection feedback.
Probe-set size and split sensitivity are treated as stability analyses rather than as part of the main selection rule.

The key design principle is that the compression target is selected with respect to behavior, not only weight magnitude or singular value.
We therefore include matched random pruning, magnitude-energy pruning, low-singular-value pruning, and uniform LoRA INT8 as controls.
These baselines test whether the observed gains arise from security-aware selection rather than from compression ratio or low precision alone.

\subsection{Materialization Operators}

We evaluate four classes of adapter materialization.
Rank pruning removes selected LoRA rank components.
Prune-then-quantize applies the same selected gate and stores the remaining adapter in low precision.
Shrink operators attenuate selected rank components rather than fully removing them.
Layer-adaptive routes assign different actions to different layer groups.
The main Qwen27B result uses the selected \method{SAC-alpha-80} route; a reported same-gate materialization applies prune-then-INT8.

\paragraph{Route labels.}
Route names are compact configuration identifiers, not separate algorithms.
Each label denotes a tuple $(\mathcal{C}, b,\lambda,T)$: component family, rank budget, behavioral scoring weights, and materialization operator.
\method{SAC-alpha-80} denotes the alpha-family gate at an 80\% rank-removal budget for the relevant component family; \method{SAC-alpha-70} and \method{SAC-alpha-60} analogously denote 70\% and 60\% budgets.
Counterfactual-score variants change only $\lambda$, the weights on TH/H/TB/B score terms.
Clean-calibrated gates include benign calibration prompts when estimating $\lambda$-weighted score drops, with the same held-out evaluation protocol used for the main rows.

\subsection{Baselines and Controls}

Hard-top-$k$ applies a hard component budget without the full security-aware routing criterion.
Matched random rank pruning uses the same reduction budget as the selected Qwen27B route but randomizes the ranking.
Magnitude-energy pruning removes components using adapter weight scale rather than behavioral evidence.
Uniform LoRA INT8 quantizes adapter weights without selective removal.
Low-singular-value rank pruning tests whether removing low-energy directions is sufficient.
AngelSlim/8-bit floating point (FP8) is reported separately as a non-\sac{} control because it evaluates external quantized model variants rather than selective adapter compression.
